% Generated by Listas/make_statements.py from pratica-lecture-05.md, item 11. Do not edit.
\textbf{Source.} Gibbons (1992), \emph{Game Theory for Applied Economists}, §2.3 Problems, Problem 2.10, p.~134 {[}PDF 141{]}. Answer in the source: none in the book. Reproduced verbatim.



\begin{quote}

2.10. The accompanying simultaneous-move game is played twice, with the outcome of the first stage observed before the second stage begins. There is no discounting. The variable \(x\) is greater than 4, so that \((4,4)\) is not an equilibrium payoff in the one-shot game. For what values of \(x\) is the following strategy (played by both players) a subgame-perfect Nash equilibrium?



\begin{quote}

Play \(Q_i\) in the first stage. If the first-stage outcome is \((Q_1, Q_2)\), play \(P_i\) in the second stage. If the first-stage outcome is \((y, Q_2)\) where \(y \neq Q_1\), play \(R_i\) in the second stage. If the first-stage outcome is \((Q_1, z)\) where \(z \neq Q_2\), play \(S_i\) in the second stage. If the first-stage outcome is \((y, z)\) where \(y \neq Q_1\) and \(z \neq Q_2\), play \(P_i\) in the second stage.

\end{quote}



{\def\LTcaptype{none} % do not increment counter

\begin{longtable}[]{@{}lllll@{}}

\toprule\noalign{}

& \(P_2\) & \(Q_2\) & \(R_2\) & \(S_2\) \\

\midrule\noalign{}

\endhead

\bottomrule\noalign{}

\endlastfoot

\(P_1\) & 2, 2 & \(x\), 0 & −1, 0 & 0, 0 \\

\(Q_1\) & 0, \(x\) & 4, 4 & −1, 0 & 0, 0 \\

\(R_1\) & 0, 0 & 0, 0 & 0, 2 & 0, 0 \\

\(S_1\) & 0, −1 & 0, −1 & −1, −1 & 2, 0 \\

\end{longtable}

}

\end{quote}

